\section{Structural Audit: Resonance Spectroscopy}

The spectral scan identifies the ``Ghost Forest'' of high-density resonances as excited baryon states. This hypothesis can be tested by experimentalists with access to hadron spectroscopy data (e.g., JLab, BESIII, LHCb).

\subsection*{Baryon Resonance Scan}

\begin{itemize}
    \item \textbf{Task:} Re-analyze the invariant mass spectra of $N^*$ and $\Delta$ resonances. Search for statistical preference for mass peaks centered at the exact geometric solutions ($D^2 + n\alpha^{-1}$ or $2D^2 - \text{offsets}$).
    \item \textbf{Expected Result:} ``Ghost'' peaks in the $E_8$ spectrum will map to low-confidence resonances in the PDG, confirming geometric quantization.
    \item \textbf{Falsification Criterion:} Lack of statistical correlation between geometric peaks and experimental resonances (Bayes Factor $< 1$).
\end{itemize}

\subsection*{Top Quark Mass Definition}

\begin{itemize}
    \item \textbf{Task:} Re-evaluate the Top Mass measurement using the $E_8$ ``Volumetric Limit'' ($337,900 \, m_e$) as a rigid constraint.
    \item \textbf{Expected Result:} Fixing the mass to the geometric limit minimizes the vacuum stability calculation uncertainty ($\lambda \approx 0$).
\end{itemize}